\section{Nakayama's lemma}
\label{section-nakayama}

\noindent
We quote from \cite{MatCA}: ``This simple but
important lemma is due to T.~Nakayama, G.~Azumaya and W.~Krull. Priority
is obscure, and although it is usually called the Lemma of Nakayama, late
Prof.~Nakayama did not like the name.''

\begin{lemma}[Nakayama's lemma]
\label{lemma-NAK}
\begin{reference}
\cite[1.M Lemma (NAK) page 11]{MatCA}
\end{reference}
\begin{history}
We quote from \cite{MatCA}: ``This simple but
important lemma is due to T.~Nakayama, G.~Azumaya and W.~Krull. Priority
is obscure, and although it is usually called the Lemma of Nakayama, late
Prof.~Nakayama did not like the name.''
\end{history}
Let $R$ be a ring with Jacobson radical $\text{rad}(R)$.
Let $M$ be an $R$-module. Let $I \subset R$
be an ideal.
\begin{enumerate}
\item
\label{item-nakayama}
If $IM = M$ and $M$ is finite, then there exists an $f \in 1 + I$ such that
$fM = 0$.
\item If $IM = M$, $M$ is finite, and $I \subset \text{rad}(R)$, then $M = 0$.
\item If $N, N' \subset M$, $M = N + IN'$, and $N'$ is finite,
then there exists an $f \in 1 + I$ such that $fM \subset N$ and $M_f = N_f$.
\item If $N, N' \subset M$, $M = N + IN'$, $N'$ is finite, and
$I \subset \text{rad}(R)$, then $M = N$.
\item If $N \to M$ is a module map, $N/IN \to M/IM$ is
surjective, and $M$ is finite, then there exists an $f \in 1 + I$
such that $N_f \to M_f$ is surjective.
\item If $N \to M$ is a module map, $N/IN \to M/IM$ is
surjective, $M$ is finite, and $I \subset \text{rad}(R)$,
then $N \to M$ is surjective.
\item If $x_1, \ldots, x_n \in M$ generate $M/IM$ and $M$ is finite,
then there exists an $f \in 1 + I$ such that $x_1, \ldots, x_n$
generate $M_f$ over $R_f$.
\item If $x_1, \ldots, x_n \in M$ generate $M/IM$, $M$ is finite, and
$I \subset \text{rad}(R)$, then $M$ is generated by $x_1, \ldots, x_n$.
\item If $IM = M$, $I$ is nilpotent, then $M = 0$.
\item If $N, N' \subset M$, $M = N + IN'$, and $I$ is nilpotent then $M = N$.
\item If $N \to M$ is a module map, $I$ is nilpotent, and $N/IN \to M/IM$
is surjective, then $N \to M$ is surjective.
\item If $\{x_\alpha\}_{\alpha \in A}$ is a set of elements of $M$
which generate $M/IM$ and $I$ is nilpotent, then $M$ is generated
by the $x_\alpha$.
\end{enumerate}
\end{lemma}

\begin{proof}
Proof of (\ref{item-nakayama}). Choose generators $y_1, \ldots, y_m$ of $M$
over $R$. For each $i$ we can write $y_i = \sum z_{ij} y_j$ with
$z_{ij} \in I$ (since $M = IM$).
In other words $\sum_j (\delta_{ij} - z_{ij})y_j = 0$.
Let $f$ be the determinant of the $m \times m$ matrix
$A = (\delta_{ij} - z_{ij})$. Note that $f \in 1 + I$
(since the matrix $A$ is entrywise congruent to the
$m \times m$ identity matrix modulo $I$).
By Lemma \ref{lemma-matrix-left-inverse} (1),
there exists an $m \times m$
matrix $B$ such that $BA = f 1_{m \times m}$. Writing out we see that
$\sum_{i} b_{hi} a_{ij} = f \delta_{hj}$ for all
$h$ and $j$; hence, $\sum_{i, j} b_{hi} a_{ij} y_j
= \sum_{j} f \delta_{hj} y_j = f y_h$ for every $h$.
In other words, $0 = f y_h$ for every $h$ (since each
$i$ satisfies $\sum_j a_{ij} y_j = 0$).
This implies that $f$ annihilates $M$.

\medskip\noindent
By Lemma \ref{lemma-contained-in-radical} an element of $1 + \text{rad}(R)$ is
invertible element of $R$. Hence we see that (\ref{item-nakayama}) implies
(2). We obtain (3) by applying (1) to $M/N$ which is finite as $N'$ is finite.
We obtain (4) by applying (2) to $M/N$ which is finite as $N'$ is finite.
We obtain (5) by applying (3) to $M$ and the submodules $\Im(N \to M)$
and $M$. We obtain (6) by applying (4) to $M$ and the submodules
$\Im(N \to M)$ and $M$.
We obtain (7) by applying (5) to the map $R^{\oplus n} \to M$,
$(a_1, \ldots, a_n) \mapsto a_1x_1 + \ldots + a_nx_n$.
We obtain (8) by applying (6) to the map $R^{\oplus n} \to M$,
$(a_1, \ldots, a_n) \mapsto a_1x_1 + \ldots + a_nx_n$.

\medskip\noindent
Part (9) holds because if $M = IM$ then $M = I^nM$ for all $n \geq 0$
and $I$ being nilpotent means $I^n = 0$ for some $n \gg 0$. Parts
(10), (11), and (12) follow from (9) by the arguments used above.
\end{proof}

\begin{lemma}
\label{lemma-NAK-localization}
Let $R$ be a ring, let $S \subset R$ be a multiplicative subset,
let $I \subset R$ be an ideal, and let $M$ be a finite $R$-module.
If $x_1, \ldots, x_r \in M$ generate $S^{-1}(M/IM)$
as an $S^{-1}(R/I)$-module, then there exists an $f \in S + I$
such that $x_1, \ldots, x_r$ generate $M_f$ as an
$R_f$-module.\footnote{Special cases: (I) $I = 0$. The lemma says
if $x_1, \ldots, x_r$ generate $S^{-1}M$, then $x_1, \ldots, x_r$
generate $M_f$ for some $f \in S$. (II) $I = \mathfrak p$ is
a prime ideal and $S = R \setminus \mathfrak p$. The lemma says if
$x_1, \ldots, x_r$ generate $M \otimes_R \kappa(\mathfrak p)$
then $x_1, \ldots, x_r$ generate $M_f$ for some
$f \in R$, $f \not \in \mathfrak p$.}
\end{lemma}

\begin{proof}
Special case $I = 0$. Let $y_1, \ldots, y_s$ be generators for $M$ over $R$.
Since $S^{-1}M$ is generated by $x_1, \ldots, x_r$, for each $i$
we can write $y_i = \sum (a_{ij}/s_{ij})x_j$ in $S^{-1}M$
for some $a_{ij} \in R$ and $s_{ij} \in S$. Multiplying by the product
$s \in S$ of the $s_{ij}$ we see that $sy_i = \sum a'_{ij}x_j$
in $S^{-1}M$ for some $a'_{ij} \in R$.
This in turn means there exist $t_i \in S$ such that
$t_isy_i = \sum t_ia'_{ij}x_j$ in $M$. Thus if $t \in S$
is the product of the $t_i$, then we
see that $y_i$ is in the $R_{st}$-submodule generated by $x_1, \ldots, x_r$
of $M_{st}$. Hence $x_1, \ldots, x_r$ generates $M_{st}$.

\medskip\noindent
General case. By the special case, we can find an $s \in S$
such that $x_1, \ldots, x_r$ generate $(M/IM)_s$ over $(R/I)_s$.
By Lemma \ref{lemma-NAK} we can find a $g \in 1 + I_s \subset R_s$
such that $x_1, \ldots, x_r$ generate $(M_s)_g$ over $(R_s)_g$.
Write $g = 1 + i/s'$. Then $f = ss' + is$ works; details omitted.
\end{proof}

\begin{lemma}
\label{lemma-when-surjective-local}
Let $A \to B$ be a local homomorphism of local rings.
Assume
\begin{enumerate}
\item $B$ is finite as an $A$-module,
\item $\mathfrak m_B$ is a finitely generated ideal,
\item $A \to B$ induces an isomorphism on residue fields, and
\item $\mathfrak m_A/\mathfrak m_A^2 \to \mathfrak m_B/\mathfrak m_B^2$
is surjective.
\end{enumerate}
Then $A \to B$ is surjective.
\end{lemma}

\begin{proof}
To show that $A \to B$ is surjective, we view it as a map of $A$-modules
and apply Lemma \ref{lemma-NAK} (6). We conclude it suffices
to show that $A/\mathfrak m_A \to B/\mathfrak m_AB$ is surjective.
As $A/\mathfrak m_A = B/\mathfrak m_B$ it suffices to show that
$\mathfrak m_AB \to \mathfrak m_B$ is surjective. View
$\mathfrak m_AB \to \mathfrak m_B$ as a map of $B$-modules and apply
Lemma \ref{lemma-NAK} (6). We conclude it suffices to see that
$\mathfrak m_AB/\mathfrak m_A\mathfrak m_B \to \mathfrak m_B/\mathfrak m_B^2$
is surjective. This follows from assumption (4).
\end{proof}
